\item We first remove the existential condition about affine lifts
from the admissibility test for a linear tuple. A finite-order
matrix $M$ admits a free finite affine lift on $T^6$ if and
only if $\ker(M-1)\ne0$. If this kernel is zero, $1-M$ is a
surjective endomorphism of the torus, so every affine map with
linear part $M$ has a fixed point. Conversely let the order of
$M$ be $d>1$ and choose a primitive $u\in\ker(M-1)\cap\mathbb Z^6$.
The map $x\mapsto Mx+u/d^2$ has actual order $d^2$: its
$d$th power is translation by $u/d$, of exact order $d$.
For each prime $p\mid d$, the generator of its order-$p$
subgroup is translation by $u/p$, which has no fixed point.
Every nontrivial stabilizer contains a prime-order subgroup,
so the whole cyclic action is free. If $M=1$, translation
by a primitive half-vector is a free lift of order two.
The invariant integral kernel is primitive, so such a
primitive vector is obtained by Smith reduction.

\item For an admissible full tuple $T=(M_1,\ldots,M_k,M_0)$ let
$H(T)$ be its maximum absolute entry. If a three-point linear
tuple has an identity at a finite position, temporarily delete
that identity, classify the remaining two-point tuple, and
reinsert the identity at the same marked position. There can
be only one such identity, since otherwise a finite matrix
would equal the inverse of the nontrivial unipotent cusp.
Deleting a fixed identity changes neither height nor the
lexicographic comparison of tuples in that conjugacy-invariant
sector. This is only a linear normalization: a nontrivial
affine order at the identity-linear site is retained in the
gluing datum. In the two-involution sector use the canonical
representative defined in \ref{maps:src:F.2.38}. This sector is invariant under simultaneous
conjugacy. In every other sector, order full integer tuples
first by $H$, then lexicographically by their row-major
entries in the displayed marked order. Define the canonical
representative to be the first admissible tuple in the
simultaneous conjugacy class.

